\section{Theory}
\label{sec:theory}

The theory answers five questions. When is a monitor valid at any stopping time (\cref{sec:theory-validity}), and when does a combination of monitors stay valid (\cref{sec:theory-combination})? How fast can any monitor grow
under a departure, and what does a mixture cost (\cref{sec:theory-growth})? What power can a monitor have at a fixed horizon (\cref{sec:theory-power})? And where do the guarantees stop: why the repeated fixed-window tests
lose their size, and why the PD monitor loses its own when the default probability of a year is random (\cref{sec:theory-failure})? Proofs are short and use only the martingale property, Ville's inequality, Jensen's inequality and
the Neyman-Pearson lemma.

\subsection{Validity at any stopping time}
\label{sec:theory-validity}

\begin{proposition}[Validity of a mixture of likelihood ratios]
\label{prop:validity}
Let the hits satisfy the null of \cref{eq:null-var}. For alternatives $p_k\in(0,1)$ and weights $w_k\ge0$, $\sum_kw_k=1$, fixed before any outcome, the process $E_t=\sum_kw_k\Lambda_t(p_k)$ of \cref{eq:mixture} is a non-negative martingale with respect to $(\mathcal F_t)$ with $E_0=1$.
For every $a\in(0,1)$,
\begin{equation}
\Prob_0\bigl(\exists\,t\ge1:\;E_t\ge 1/a\bigr)\le a ,
\label{eq:ville}
\end{equation}
and for every stopping time $\sigma$ the quantity $\min\bigl(1,1/\max_{s\le\sigma}E_s\bigr)$ is a p-value: its probability of falling at or below $u$ is at most $u$.
\end{proposition}
\begin{proof}
Given $\mathcal F_{t-1}$ the factor of day $t$ in $\Lambda_t(p)$ takes the value $p/\alpha$ with probability $\alpha$ and $(1-p)/(1-\alpha)$ with probability $1-\alpha$, so its conditional mean is $p+(1-p)=1$. Hence $\Lambda_t(p)$ is a
martingale with $\Lambda_0(p)=1$, and a convex combination of martingales with weights that do not depend on the data is a martingale. Let $\sigma=\inf\{t:E_t\ge1/a\}$ and $N$ a finite integer. The stopping time $\sigma\wedge N$ is bounded, so optional stopping gives
$1=\E[E_{\sigma\wedge N}]\ge\E\bigl[E_\sigma\indic\{\sigma\le N\}\bigr]\ge a^{-1}\Prob_0(\sigma\le N)$. Letting $N\to\infty$ gives \cref{eq:ville}. The event $\{\max_{s\le\sigma}E_s\ge1/u\}$ is contained in $\{\exists t:E_t\ge1/u\}$, whose probability is at most $u$.
\end{proof}

The proposition needs the conditional null only, not independence: the hit probability given the past is $\alpha$. It holds for any weights and any grid, so a poor choice of alternatives costs power and never size.
It also holds for a monitor that a person stops by hand at a time chosen after looking at the path, which is the sense of optional stopping. The anytime-valid p-value of \cref{eq:rule} is the p-value of the proposition at the stopping time that is the present day.

\begin{proposition}[Validity of the clustering and PD monitors]
\label{prop:clustering}
Under the null of \cref{eq:null-var}, $E^{\mathrm{cl}}_t$ of \cref{eq:cluster} is a non-negative martingale with $E_0^{\mathrm{cl}}=1$. Under the null of \cref{eq:null-pd}, $E_{t,g}$ of \cref{eq:pd-lr} is a non-negative martingale with $E_{0,g}=1$ for every grade, and \cref{eq:ville} holds for both.
\end{proposition}
\begin{proof}
For the clustering factor, a day that follows a non-hit contributes one. If yesterday was a hit, the factor of today is $q/\alpha$ with probability $\alpha$ and $(1-q)/(1-\alpha)$ with probability $1-\alpha$, because the null fixes the conditional probability of a hit at $\alpha$ whatever the past; the mean is one. For the PD factor, with $r=m(1-p)/(1-mp)$ the factor of a step is $\bigl((1-mp)/(1-p)\bigr)^{n}r^{d}$, and for $d\sim\mathrm{Binomial}(n,p)$ one has $\E[r^d]=(1-p+pr)^n=\bigl((1-p)/(1-mp)\bigr)^n$, so the mean of the factor is one. Mixtures with fixed weights preserve the martingale property, and \cref{eq:ville} follows as in \cref{prop:validity}.
\end{proof}

\subsection{Combination rules}
\label{sec:theory-combination}

\begin{proposition}[Combination of e-processes]
\label{prop:product}
Let $E^1_t,\dots,E^G_t$ be e-processes. Their average is an e-process whatever their dependence. Their product is an e-process if, given $\mathcal F_{t-1}$, the factors of step $t$ are independent. The product of two e-processes that are functions of the same
observation is an e-process only if each is constant: if $E^1=E^2=E$ and $E$ has mean one, then $\E[E^1E^2]=1+\Var(E)$.
\end{proposition}
\begin{proof}
Expectation is linear, so $\E[\tfrac1G\sum_gE^g_t\mid\mathcal F_{t-1}]\le\tfrac1G\sum_gE^g_{t-1}$. For independent factors, the conditional expectation of the product is the product of the conditional expectations, each at most one. If $E^1=E^2=E$ then $\E[E^2]=\Var(E)+\E[E]^2=1+\Var(E)$.
\end{proof}

The third statement is the reason the battery of the two VaR monitors averages. Two computations in the experiments show its size. For two grades with identical defaults (perfect dependence), three obligors per grade and a PD of $0.3$, the exact expectation of the
product e-process after one step is \nProductPdProduct\ and that of the average is \nProductPdAverage. For the rate and clustering monitors read on one stream of $10$ days at $\alpha=0.2$, the exact expectation of their product reaches \nProductMax\ (\cref{tab:null-expectation}), so a product on one stream would reject almost every correct model.
The enumeration of all \nEnumSequences\ hit sequences of $12$ days confirms the martingale property of the valid processes: the largest departure of their expectation from one is $\nEnumWorst$, floating-point error.

\begin{table}[htbp]
\centering
\small
\caption{Exact null expectation of each process by enumeration of all 4,096 hit sequences of 12 days at a VaR level of 0.2: the largest departure of the expectation from 1 over the 12 days is at the level of floating-point error for every valid process, and the last row shows the largest expectation of the product of two e-processes read on the same stream over ten days, which is not an e-process.}
\label{tab:null-expectation}
\resizebox{\linewidth}{!}{%
\begin{tabular}{p{7.4cm}rr}
\toprule
Process & Largest departure of $\mathbb E[E_t]$ from 1 over $t\le 12$ & $\mathbb E[E_{12}]$ \\
\midrule
\relax Rate e-process, two-sided & \relax 1.2e-14 & 1.000000000000 \\
\relax Rate e-process, upper & \relax 9.8e-15 & 1.000000000000 \\
\relax Rate e-process, lower & \relax 2.9e-14 & 1.000000000000 \\
\relax Clustering e-process & \relax 3.6e-15 & 1.000000000000 \\
\relax Average of the rate and clustering e-processes & \relax 1.0e-14 & 1.000000000000 \\
\relax E-backtest of Wang, Wang and Ziegel (2022) & \relax 4.7e-15 & 1.000000000000 \\
\relax Uniform-prior Bayes factor & \relax 5.1e-15 & 1.000000000000 \\
\relax Product of the rate and clustering e-processes on one stream ($T=10$, $\alpha=0.2$) & \relax 494 & -- \\
\bottomrule
\end{tabular}}
\end{table}


\subsection{Growth rate and delay}
\label{sec:theory-growth}

\begin{proposition}[Growth rate of an e-process]
\label{prop:growth}
Let the hits be independent Bernoulli($p$) with $p\ne\alpha$ and write $\mathrm{KL}(p\Vert r)=p\log\frac pr+(1-p)\log\frac{1-p}{1-r}$. Then (i) $\E_p[\log\Lambda_t(p_k)]=t\bigl[\mathrm{KL}(p\Vert\alpha)-\mathrm{KL}(p\Vert p_k)\bigr]$. (ii) For the equal-weight mixture of $K$ alternatives,
\begin{equation}
\E_p[\log E_t]\;\ge\;t\max_k\bigl[\mathrm{KL}(p\Vert\alpha)-\mathrm{KL}(p\Vert p_k)\bigr]-\log K .
\label{eq:growth}
\end{equation}
(iii) Every non-negative process $E$ with $\E_0[E_t]\le1$ that is a function of the first $t$ hits satisfies $\E_p[\log E_t]\le t\,\mathrm{KL}(p\Vert\alpha)$, with equality for $\Lambda_t(p)$.
\end{proposition}
\begin{proof}
(i) Take the expectation of the daily term $I\log(p_k/\alpha)+(1-I)\log\frac{1-p_k}{1-\alpha}$ and write it as the difference of two divergences. (ii) A mixture exceeds each of its terms, so $\log E_t\ge\log w_j+\log\Lambda_t(p_j)$ for every $j$; use (i) with $w_j=1/K$.
(iii) Let $P_0^{(t)}$ and $P_p^{(t)}$ be the laws of $t$ hits. Then $\E_p[\log E_t]-t\,\mathrm{KL}(p\Vert\alpha)=\E_p\bigl[\log\bigl(E_t\,dP_0^{(t)}/dP_p^{(t)}\bigr)\bigr]\le\log\E_0[E_t]\le0$ by Jensen's inequality.
\end{proof}

Part (iii) says no e-process grows faster in expectation than the likelihood ratio of the true alternative, and (ii) says the mixture loses at most $\log K$ in the exponent, plus the mismatch $\mathrm{KL}(p\Vert p_k)$ of the closest alternative.
Together they give the first-order delay: the expected log of the process reaches $\log(1/a)$ after about $\log(1/a)/\mathrm{KL}(p\Vert\alpha)$ days for the matched likelihood ratio and after at most $(\log(1/a)+\log K)/g^\ast$ days for the mixture, where $g^\ast$ is the best growth rate of \cref{eq:growth}.
At $\alpha=0.01$ and a hit rate three times the nominal, $\mathrm{KL}(0.03\Vert0.01)=\nKlThree\times10^{-3}$ nats per day, so the matched bet needs about \nKlSprtThree\ days and the mixture of eight alternatives about \nKlMixThree\ days; at four times the nominal they need \nKlSprtFour\ and \nKlMixFour\ days
(\cref{tab:theory-delay}). These are first-order times for the expected log; the median first passage differs from them because a hit adds a jump to the path. The measured medians of \cref{sec:results} lie below the first-order times: \nMedThreeSprt\ days against \nKlSprtThree\ for the matched bet and \nMedThreeE\ days against \nKlMixThree\ for the mixture at three times the nominal rate, and they keep the order the theory predicts.

\subsection{Power at a fixed horizon}
\label{sec:theory-power}

\begin{proposition}[Power of a monitor at a fixed horizon]
\label{prop:power}
Let a monitor have a size $s_n=\Prob_0(\tau\le n)$ by day $n$ and let $\beta_n(s)$ be the power at level $s$ of the most powerful randomised test of the null against the alternative on the first $n$ hits, as the Neyman-Pearson lemma defines it. Then $\Prob_1(\tau\le n)\le\beta_n(s_n)$.
\end{proposition}
\begin{proof}
The event $\{\tau\le n\}$ is determined by the first $n$ hits, so it is a randomisation-free test of size $s_n$ of the null against the alternative. The power of every such test is at most that of the most powerful one.
\end{proof}

The proposition is the structural reason for the power cost. A monitor that is valid at every time has spent only part of its size by day $n$: at $\alpha=0.01$ and $n=250$ the rate e-process has rejected \nNullByE\,\% of the correct series by day 250, so no monitor with that spending can have more than \nExactThreeNpE\,\% power at a hit rate three times the nominal, against
\nExactThreeKupUp\,\% for the upper tail of the Kupiec test, which spends \nExactKupUpperSize\,\% of its size there (\cref{tab:size-matched}). The gap between the two powers is the cost of validity at later days, and a monitor can only narrow it by spending more of its size early, at the expense of the later days.

\subsection{Where the guarantees stop}
\label{sec:theory-failure}

\begin{proposition}[Repeated fixed-window tests have size one in the limit]
\label{prop:repeated}
Let $S_t$ be the number of hits in $t$ days under the null and $Z_t=(S_t-t\alpha)/\sqrt{t\alpha(1-\alpha)}$. For every $c>0$, $\Prob_0(\sup_t|Z_t|\ge c)=1$, so the Kupiec test read on all days so far at any fixed level has a size that tends to one with the horizon. A test of a window of fixed length $w$ with one-shot size $s_w>0$, read on every day, rejects a correct model by day $T$ with probability at least
$1-(1-s_w)^{\lfloor T/w\rfloor}$.
\end{proposition}
\begin{proof}
By the law of the iterated logarithm, $\limsup_tZ_t/\sqrt{2\log\log t}=1$ almost surely, so $|Z_t|$ exceeds any constant infinitely often; the Kupiec statistic equals $Z_t^2$ up to a vanishing term. For a window of length $w$, the non-overlapping windows $[(j-1)w+1,jw]$ are independent and each is rejected with probability $s_w$, so the probability that none of $\lfloor T/w\rfloor$ of them rejects is $(1-s_w)^{\lfloor T/w\rfloor}$.
\end{proof}

The proposition explains the growth of the repeated-test curves of \cref{fig:size-horizon}: the Kupiec test on the last 250 days and the traffic light, whose one-shot sizes are \nHorizonWindowKupTrail\,\% and \nHorizonWindowLight\,\% by day 250, reject \nHorizonEndKupTrail\,\% and \nHorizonEndLight\,\% of correct series by day \nHorizonDays.
It also shows that a boundary calibrated by simulation to a size of 5\,\% over 1,000 days is not a guarantee beyond 1,000 days.

\begin{proposition}[Loss of validity of the PD monitor when the default probability is random]
\label{prop:overdispersion}
Let $d\mid P\sim\mathrm{Binomial}(n,P)$ with $P\in(0,1)$ random, $\E[P]=p$, $n\ge2$ and $P$ not almost surely constant, and let $m\ne1$ with $mp<1$. The one-step factor $f_m(d)=\bigl((1-mp)/(1-p)\bigr)^{n}r^{d}$ with $r=m(1-p)/(1-mp)$ satisfies $\E[f_m(d)]>1$. Hence a mixture $\sum_kw_k\prod_sf_{m_k}$ over steps that are independent draws of this kind has a conditional mean above one at every step, and
$\E[E_t]$ grows geometrically with $t$.
\end{proposition}
\begin{proof}
Given $P$, $\E[r^d\mid P]=(1+(r-1)P)^n$, so $\E[f_m(d)]=c\,\E\bigl[(1+(r-1)P)^n\bigr]$ with $c=((1-mp)/(1-p))^n$. The function $x\mapsto(1+(r-1)x)^n$ is strictly convex on the interval where it is positive when $n\ge2$ and $r\ne1$. By Jensen's inequality its expectation exceeds its value at $\E[P]=p$, which is $c^{-1}$ by the computation in the proof of \cref{prop:clustering}. Therefore $\E[f_m]>1$ for every $m\ne1$ and the same holds for any mixture.
\end{proof}

The proposition is the theoretical form of the assumption stated in \cref{sec:problem}: the PD monitor is valid when the defaults of a step are conditionally binomial, and a random default probability, which is what an asset correlation produces in the one-factor model, makes every factor exceed its target in expectation, so the process drifts upwards under a correct PD.
The experiment of \cref{sec:results-pd} measures the effect: the size is far above 5\,\% already at an asset correlation of $0.03$.
